% SOURCE_AUTHORITY: sga5_fr_workpass.tex, lines 11455--11516 (LF-normalized unit).
% SOURCE_SHA256: 791F4EFFC5E02832D5D77ED03518C8156D6F07E4C8238B03545DB93D883FBB28
\emph{Demonstração.} A última afirmação não é senão (9.1.1). Do mesmo modo, (9.6.2) exprime que $\mu\circ\lambda=0$. Demonstremos que $\mu$ é sobrejetiva. Para isso precisaremos de
\[
f_*f^*(z)=z\qquad(z\in C^{\bullet}(X)),
\tag{9.9.1}
\]
que resulta imediatamente da fórmula de projeção para $f$ e do fato evidente de que $f_*(1)=1$. Seja, então, $t\in C^{\bullet}(X')$. Por (9.9.1), temos
\[
f_*(t-f^*f_*(t))=0.
\]
Como $f$ induz um isomorfismo de $U'=X'-Y'$ sobre $U=X-Y$, deduz-se que a restrição a $U'$ de $t-f^*f_*(t)$ é nula e, por conseguinte ([CH] 4 th. 2), tem-se uma igualdade
\[
t-f^*f_*(t)=j_*(a)
\qquad\text{com }a\in C^{\bullet}(Y').
\]
Demonstremos agora a exatidão no termo central. Sejam, pois, $a\in C^{p-1}(Y')$ e $b\in C^p(X)$ tais que
\[
j_*(a)+f^*(b)=0.
\tag{9.9.2}
\]
Aplicando $f_*$ aos dois membros de (9.9.2), vemos que
\[
f_*j_*(a)+f_*f^*(b)=0,
\]
de onde, tendo em conta (9.9.1),
\[
b=-i_*g_*(a).
\tag{9.9.3}
\]
Colocando $x=g_*(a)$, basta, então, provar que
\[
a=g^*(x)c_{d-1}(\check F).
\tag{9.9.4}
\]
Ponhamos, para isso,
\[
z=a-g^*(x)c_{d-1}(\check F)\qquad(x=g_*(a)),
\]
e mostremos que $z=0$. Para fazer isso, veremos que temos as duas relações $j_*(z)=0$ e $g_*(z)=0$, e depois que todo elemento $z$ que verifique ambas as relações é necessariamente nulo. Por definição,
\[
j_*(z)=j_*(a)-j_*(g^*(x)\cdot c_{d-1}(\check F)),
\]
Isto é, tendo em conta (9.6.2),
\[
j_*(z)=j_*(a)-f^*i_*(x).
\]
Substituindo $i_*(x)$ por $-b$ (9.9.3) na expressão anterior, deduz-se de (9.9.2) que $j_*(z)=0$. Considerando a definição de $x$, vemos que
\[
g_*(z)=x-g_*(g^*(x)c_{d-1}(\check F))=0
\]
pela última afirmação do teorema. Agora vamos demonstrar como a nulidade de $z$ resulta das duas relações $j_*(z)=0$ e $g_*(z)=0$. Como $Y'\simeq\mathbb P_Y(N)$, o elemento $z$ admite uma representação única da forma
\[
z=g^*(a_0)+g^*(a_1)\xi+\cdots+g^*(a_{d-1})\xi^{d-1},
\]
com $a_i\in C^{p-1-i}(Y)$. Como $g_*(z)=0$, resulta da fórmula de projeção para $g$ e do fato de que $g_*$ diminui os graus em $d-1$ unidades que $a_{d-1}=0$. Por outro lado, por (9.2),
\[
z\,\xi=j^*j_*(z)=0,
\]
de onde
\[
g^*(a_0)\xi+\cdots+g^*(a_{d-2})\xi^{d-1}=0;
\]
Assim, todos os $a_i$ são nulos e $z=0$.
